\section{Introduction}
\label{sec:intro}

\subsection{The gap between a competent model and a competent operator}

The current wave of industrial AI rests on a simple observation: a large
language model that can read a manual, write a query and call a function is,
in principle, a general-purpose operator's assistant. The demonstrations are
compelling. Summarise the shift log, explain an alarm, draft the structured
text for a new interlock, propose a setpoint. The deployments are not. The
failure is rarely a failure of fluency or even of reasoning; it is a failure of
\emph{reference}. The model does not know which physical thing the words denote.

Consider a request that any operator answers in three seconds: \emph{``what is
the outlet pressure of the Line~1 filler?''} In the reference plant used
throughout this paper, five distinct tags carry the description string
``Outlet pressure''. Two of them are on fillers, on different production lines,
commissioned nine years apart; one is calibrated in \texttt{bar} and the other
in \texttt{psi}. A sixth tag, on the same filler, is the pressure
\emph{setpoint} rather than the measurement, and is writable where the other is
not. A model handed a historian export sees six rows of near-identical English
and a column of numbers. It will answer. It may well answer correctly, because
the tag mnemonics encode the plant hierarchy for a reader who knows the
convention. But nothing in the representation makes the answer \emph{checkable},
and nothing prevents the same model, in the same session, from proposing a write
to the read-only measurement instead of the setpoint, in the wrong unit, while
the machine is in a state that forbids the change.

This is the problem that semantic technologies address, and industrial
automation has been addressing it, at considerable expense, for three decades:
\opcua{} information models, the Asset Administration Shell, ECLASS property
dictionaries, AutomationML engineering exchange, ISA-95 equipment hierarchies,
PackML state machines and the Module Type Package. From the web side, the list
includes RDF, OWL, SHACL, SPARQL, QUDT, SOSA/SSN and a shelf of domain ontologies. The
motivation was machine-to-machine interoperability, not machine comprehension.
The question this paper asks is whether that investment now pays a second
dividend: \emph{how much of what a language agent needs is already standardised,
which pieces are missing, and what does each piece actually buy in measurable
accuracy?}

\subsection{Why this is a survey \emph{and} an experiment}

Surveys of industrial semantic technology exist, and so does a fast-growing
literature on grounding language models with knowledge graphs. They do not meet.
Standards surveys evaluate technologies against interoperability requirements
such as coverage, tool support and conformance, which are the wrong criteria for
an agent. Knowledge-graph--RAG papers evaluate on open-domain question answering
over encyclopaedic graphs, where there are no units, no interlocks, no
read-only variables and no consequence to being wrong.

We therefore do two things at once. First (\S\ref{sec:layer0}--\S\ref{sec:layer4})
we survey the field through a lens built for agents: a ten-criterion
\emph{agent-usability framework} (\S\ref{sec:framework}) that asks of every
technology what it lets an agent \emph{resolve}, \emph{type}, \emph{convert},
\emph{traverse}, \emph{do}, \emph{check}, \emph{trust}, \emph{query},
\emph{afford} and \emph{deploy}. Second (\S\ref{sec:design}--\S\ref{sec:results})
we build a reproducible testbed in which those criteria are ablated one layer at
a time, and measure the consequences.

The experimental design deliberately avoids benchmarking a particular model.
Model rankings expire within months, and an accuracy delta measured on one model
tells you little about the representation, because the reader is not held
constant. Instead we measure properties of the representation itself:

\begin{itemize}
  \item \textbf{Grounding.} Can a fixed resolution algorithm map a
  natural-language mention onto the right tag, given only what this
  representation exposes?
  \item \textbf{Sufficiency ceiling.} Does the retrieved context contain
  \emph{every} fact the question requires? A perfect reader cannot exceed this;
  a gap here is a gap no prompting, fine-tuning or scaling can close.
  \item \textbf{Guardrail efficacy.} Of a corpus of faulty agent actions,
  how many can be \emph{proved} wrong, and how many legal actions are wrongly
  rejected?
  \item \textbf{Action-space fidelity.} How much of the emittable call space
  does the representation eliminate, and does it eliminate the right part?
  \item \textbf{Context economy.} What does all this cost in tokens, when is
  it cheaper to run a query than to fill a window, and does a rich
  representation waste its budget on layers the question never asked about?
  \item \textbf{Robustness.} How much of the un-grounded baseline's apparent
  competence is really the tag naming convention, and what happens in a plant
  whose convention is different, absent, or untruthful?
\end{itemize}

\subsection{Contributions}

\begin{enumerate}[label=\textbf{C\arabic*.},leftmargin=2.6em]
  \item \textbf{A layered map of industrial semantics for agents.} We organise
  42 standards and vocabularies into five layers
  (L0 identification, L1 information models, L2 formal semantics, L3 domain
  ontologies, L4 dataspaces) and show that the layers are not interchangeable:
  each contributes a distinct and non-substitutable agent affordance
  (\S\ref{sec:framework}--\S\ref{sec:layer4}).

  \item \textbf{The agent-usability framework (AUF).} Ten criteria with a
  published four-point rubric (Appendix~\ref{app:rubric}), applied uniformly to
  every surveyed technology, yielding a comparison matrix
  (Table~\ref{tab:auf-matrix}) that makes the field's blind spots visible: unit
  semantics, behavioural semantics and constraint semantics are supported by far
  fewer standards than type semantics.

  \item \textbf{A reproducible industrial grounding benchmark.} A reference
  bottling plant with \NAssets{} assets and \NSignals{} signals, projected into
  seven interchange formats (RDF, \opcua{} NodeSet2, AAS~JSON, WoT~TD, NGSI-LD,
  JSON-LD and a historian CSV) from one declarative source, re-taggable under
  \NSchemes{} naming conventions, plus \NTasks{} agent tasks with declared
  evidence requirements and \NCases{} labelled proposed agent actions across
  twelve fault families (\S\ref{sec:design}).

  \item \textbf{Nine experiments quantifying what each semantic layer buys.}
  Grounding, sufficiency, serialisation economy, guardrail efficacy,
  action-space reduction, the accuracy ceiling, query-versus-stuffing,
  robustness to the tag naming convention, and the effect of intent-routed
  projection (\S\ref{sec:results}).

  \item \textbf{Ten findings and a design pattern.} Chief among them: a fully
  typed information model, which is the current state of practice, is a
  \emph{local optimum} that leaves the most dangerous failure modes undetectable
  (\S\ref{sec:discussion}); and constraint checking should be treated as a
  first-class deliverable of the semantic layer rather than a by-product.

  \item \textbf{An open artefact.} Every table, figure and number in this paper
  is regenerated by \texttt{python run\_experiments.py}, offline, deterministically,
  in under a minute (Appendix~\ref{app:artifacts}).
\end{enumerate}

\subsection{The thesis, in one paragraph}

Semantics is not one thing, and the parts of it that industry has standardised
best are not the parts that agents need most. Identification and typing, the
mature part of the field, lift entity grounding from \GrndTopOneCZero{} to
\GrndTopOneCTwo{} and let a validator catch bad references, bad ranges and bad
permissions. They do essentially nothing for the failures that hurt: a command
that is illegal in the current machine state, a write that violates an
interlock, a number that is right in the wrong unit, an answer that contradicts
the observation it claims to quote. Those require unit ontologies, behavioural
models, explicit constraint models and provenance, and adding them lifts
guardrail recall from \RecallGTwo{} to \RecallGFour{} and the task ceiling from
\CeilingCTwo{} to \CeilingCFour{}. Two corollaries matter commercially: solving
grounding by shipping a bigger, better-typed catalogue into a bigger context
window is measurably a dead end (Figure~\ref{fig:ceiling}b); and the competence
that makes an un-grounded pipeline look adequate is largely the tag naming
convention, which is semantics that cannot be audited (\S\ref{sec:e8}).
